% Hafta 12 -- alıştırma gövdesi (ders notu ve ayrı alıştırma kağıdı tarafından \input edilir)
\section{Çözümlü Alıştırmalar}
Her soruyu önce kendiniz deneyin, sonra çözümü okuyun.

\begin{alistirma}
\(\dfrac{a}{b}=\dfrac{4}{7}\) ve \(a+b=44\) olduğuna göre \(b-a\) kaçtır?
\end{alistirma}
\begin{cozum}
\(a=4k\), \(b=7k\) yazalım: \ \(4k+7k=11k=44\Rightarrow k=4\).\\
\(a=16\), \(b=28\) \(\Rightarrow\) \(b-a=12\) \ (kısa yol: \(b-a=3k=12\)).
\end{cozum}

\begin{alistirma}
\(\dfrac{x}{4}=\dfrac{y}{5}=\dfrac{z}{6}\) ve \(x+y+z=45\) olduğuna göre \(x\cdot z\) kaçtır?
\end{alistirma}
\begin{cozum}
\(x=4k,\ y=5k,\ z=6k\): \ \(15k=45\Rightarrow k=3\).\\
\(x=12\), \(z=18\) \(\Rightarrow\) \(x\cdot z=216\).
\end{cozum}

\begin{alistirma}
\(2a=5b\) ve \(a-b=18\) olduğuna göre \(a\) ve \(b\) kaçtır?
\end{alistirma}
\begin{cozum}
\(2a=5b=10k\) diyelim (\(\mathrm{EKOK}(2,5)=10\)): \ \(a=5k\), \(b=2k\).\\
\(a-b=3k=18\Rightarrow k=6\) \(\Rightarrow\) \(a=30\), \(b=12\). \ Kontrol: \(2\cdot 30=60=5\cdot 12\) \checkmark
\end{cozum}

\begin{alistirma}
\(2600\) TL, üç kardeş arasında \(\dfrac{1}{2}\), \(\dfrac{1}{3}\) ve \(\dfrac{1}{4}\) sayıları ile orantılı olarak paylaştırılıyor. Her kardeş kaç TL alır?
\end{alistirma}
\begin{cozum}
Paydaların EKOK'u \(12\) ile çarpalım: \(\tfrac12:\tfrac13:\tfrac14=6:4:3\).\\
\(6k+4k+3k=13k=2600\Rightarrow k=200\) \(\Rightarrow\) paylar \(1200\), \(800\), \(600\) TL.\\
(Bu, \(2,3,4\) ile ters orantılı paylaştırma ile aynıdır.)
\end{cozum}

\begin{alistirma}
\(9\) işçi bir işi \(20\) günde bitirebiliyor. \(5\) gün çalıştıktan sonra aynı nitelikte \(6\) işçi daha katılıyor. Kalan iş kaç günde biter?
\end{alistirma}
\begin{cozum}
\(9\) işçi \(5\) gün çalıştı; kalan iş \(9\) işçinin \(20-5=15\) günlük işidir.\\
Kalan iş için işçi sayısı ile gün ters orantılı: \(9\cdot 15=(9+6)\cdot x\Rightarrow 135=15x\Rightarrow x=9\) gün.
\end{cozum}

\begin{alistirma}
Özdeş \(3\) makine günde \(8\) saat çalışarak \(6\) günde \(1800\) parça üretiyor. Bu makinelerden \(4\) tanesi günde \(6\) saat çalışarak \(2400\) parçayı kaç günde üretir?
\end{alistirma}
\begin{cozum}
\[\frac{1800}{3\cdot 8\cdot 6}=\frac{2400}{4\cdot 6\cdot x}\ \Rightarrow\ \frac{1800}{144}=\frac{2400}{24x}\ \Rightarrow\ \frac{25}{2}=\frac{100}{x}\ \Rightarrow\ x=8\text{ gün}\]
(Bir makine bir saatte \(1800\div 144=12{,}5\) parça üretir; \(4\cdot 6\cdot 12{,}5=300\) parça/gün, \(2400\div 300=8\).)
\end{cozum}

\begin{alistirma}
\(25\) öğrencilik bir sınıfın not ortalaması \(64\) olarak hesaplanmıştır. Daha sonra bir öğrencinin \(82\) olan notunun yanlışlıkla \(32\) yazıldığı fark ediliyor. Sınıfın doğru ortalaması kaçtır?
\end{alistirma}
\begin{cozum}
Hatalı toplam: \(25\cdot 64=1600\). Doğru toplam: \(1600-32+82=1650\).\\
Doğru ortalama: \(\dfrac{1650}{25}=66\) \ (kısa yol: toplam \(50\) arttı, ortalama \(\frac{50}{25}=2\) artar).
\end{cozum}

\begin{alistirma}
İki pozitif sayının aritmetik ortalaması \(10\), geometrik ortalaması \(8\)'dir. Bu sayılar ve harmonik ortalamaları kaçtır?
\end{alistirma}
\begin{cozum}
\(a+b=20\), \(ab=64\): sayılar \(t^2-20t+64=(t-4)(t-16)=0\) kökleri, yani \(4\) ve \(16\).
\(\text{HO}=\frac{2ab}{a+b}=\frac{128}{20}=6{,}4\); \ \(10\ge 8\ge 6{,}4\) \checkmark
\end{cozum}

\kendinizoncesi
\section{Alıştırmalar (Kendiniz Çözün)}

\begin{enumerate}[label=\textbf{\arabic*.},leftmargin=*,itemsep=2pt]
\item \(\dfrac{x+5}{x-1}=\dfrac{7}{4}\) olduğuna göre \(x\) kaçtır?
\item \(\dfrac{a}{b}=\dfrac{3}{5}\) olduğuna göre \(\dfrac{a+b}{b-a}\) kaçtır?
\item \(\dfrac{a}{5}=\dfrac{b}{6}=\dfrac{c}{9}\) ve \(a+b+c=60\) olduğuna göre \(c\) kaçtır?
\item \(720\) TL, \(3\), \(4\) ve \(5\) ile doğru orantılı paylaştırılıyor. En az pay kaç TL'dir?
\item \(1:300\,000\) ölçekli bir haritada \(8\) cm olan uzaklık gerçekte kaç km'dir?
\item Birbirini çeviren iki dişli çarktan birinde \(40\), diğerinde \(60\) diş vardır. Küçük çark \(12\) tur attığında büyük çark kaç tur atar?
\item Eşit kapasiteli \(6\) musluk boş bir havuzu \(20\) saatte dolduruyor. Musluk sayısı \(2\) katına çıkarılırsa havuz kaç saatte dolar?
\item \(x\); \(y\) ile doğru, \(z\) ile ters orantılıdır. \(y=3\), \(z=4\) iken \(x=6\)'dır. \(y=5\), \(z=10\) iken \(x\) kaçtır?
\item \(5\) usta günde \(8\) saat çalışarak \(6\) günde \(120\) çift ayakkabı yapıyor. \(4\) usta günde \(10\) saat çalışarak \(160\) çift ayakkabıyı kaç günde yapar?
\item \(20\) kişilik bir grubun ortalaması \(55\), \(30\) kişilik bir grubun ortalaması \(65\)'tir. İki grubun birlikte ortalaması kaçtır?
\item \(5\) ve \(20\) sayılarının geometrik ve harmonik ortalamaları kaçtır?
\item Bir araç bir yolu saatte \(30\) km hızla gidip aynı yolu saatte \(70\) km hızla dönüyor. Ortalama hızı kaç km/sa'tir?
\item Ardışık \(11\) tek sayının toplamı \(231\)'dir. En büyüğü kaçtır?
\item \(a\) ile \(b\)'nin aritmetik ortalaması \(13\), geometrik ortalaması \(12\)'dir. \(a^2+b^2\) kaçtır?
\end{enumerate}

\needspace{8\baselineskip}
\subsection*{Cevap Anahtarı}
{\small
\begin{multicols}{2}
\begin{enumerate}[label=\textbf{\arabic*.},leftmargin=*,itemsep=1pt]
\item \(4x+20=7x-7\Rightarrow x=9\).
\item \(a=3k,\ b=5k\): \(\frac{8k}{2k}=4\).
\item \(20k=60\Rightarrow k=3\); \(c=27\).
\item \(12k=720\Rightarrow k=60\); \(3\cdot 60=180\) TL.
\item \(8\cdot 300\,000=2\,400\,000\) cm \(=24\) km.
\item \(40\cdot 12=60\cdot x\Rightarrow x=8\) tur.
\item \(6\cdot 20=12\cdot t\Rightarrow t=10\) saat.
\item \(k=\frac{6\cdot 4}{3}=8\); \(x=8\cdot\frac{5}{10}=4\).
\item \(\frac{120}{5\cdot 8\cdot 6}=\frac{160}{4\cdot 10\cdot x}\Rightarrow\frac12=\frac{4}{x}\Rightarrow x=8\) gün.
\item \(\frac{1100+1950}{50}=\frac{3050}{50}=61\).
\item \(\text{GO}=\sqrt{100}=10\); \(\text{HO}=\frac{200}{25}=8\).
\item \(\frac{2\cdot 30\cdot 70}{100}=42\) km/sa.
\item Ortanca \(\frac{231}{11}=21\); en büyük \(21+10=31\).
\item \(a+b=26,\ ab=144\): \(676-288=388\).
\end{enumerate}
\end{multicols}}
